\documentclass[10pt,a4paper]{amsart}
\usepackage[margin=19mm]{geometry}
\usepackage{amsmath,amssymb,mathtools}
\usepackage[scr=rsfs,cal=euler]{mathalfa}
\usepackage{xfrac}
\usepackage[unicode,hidelinks,bookmarks=false]{hyperref}
\newcommand{\ie}{i.e.\ }
\newcommand{\cf}{cf.\ }
\newcommand{\resp}{respectively\ }
\newcommand{\Subsection}{\subsection}
\providecommand{\nobreakdash}{\nobreak\hbox{-}}
\setlength{\emergencystretch}{2em}
\multlinegap0pt
\usepackage{float}
\usepackage{amssymb}
\usepackage{xcolor}
\usepackage{bm}
\usepackage{tikz}
\newtheorem{corollary}{Corollary}
\newtheorem{proposition}{Proposition}
\DeclareMathOperator{\Aut}{Aut}
\newcommand{\GG}{\mathcal{G}}
\newcommand{\QQ}{\mathbb{Q}}
\newcommand{\RR}{\mathbb{R}}
\newcommand{\ZZ}{\mathbb{Z}}
\newcommand{\bb}[1]{\boldsymbol{#1}}
\renewcommand{\d}{\mathrm{d}}
\title{\bf Lee--Yang phenomena in \\edge-coloured graph counting}
\author{Maximilian Wiesmann}
\date{Source: arXiv:2601.02525v1 (CC BY 4.0)}
\begin{document}
\maketitle

\noindent\emph{Source and licence.} \bf Lee--Yang phenomena in \\edge-coloured graph counting. Version arXiv:2601.02525v1 (CC BY 4.0), distributed by arXiv under the Creative Commons Attribution 4.0 International licence, http://creativecommons.org/licenses/by/4.0/, as recorded in the arXiv licence metadata for this exact version and shown on the abstract page https://arxiv.org/abs/2601.02525v1. Full licence notice: \texttt{licenses/arxiv-2601.02525-CC-BY-4.0.txt}.\par\smallskip

\noindent\textit{Editorial scope.} Counted unit: the article's proposition prop:An\_exp\_integral (source0.tex lines 470--478). The article's complete original proof of it is delivered with it (source0.tex lines 480--511, one \texttt{proof} environment of 32 lines). No mathematics is written, repaired or re-derived: the statement and the proof are the article's own lines, in the article's own notation, and no retained line is substituted or re-flowed.  Retained uncounted as the source-local context the proof needs: lines 350--357.  Nothing was omitted from any retained span, so the package carries no omission record.  No retained line contains a citation, so the wrapper carries no bibliography and no bibliography entry.  No retained line contains a figure environment, an image-inclusion command or drawing code, so no image is delivered and the package declares no asset dependency.  Editorial formatting only, in addition: an AMS \texttt{amsart} preamble that restates the article's own declarations and macros used by the retained lines, namely the environments \texttt{corollary}, \texttt{proposition} on separate counters, together with the standard packages those lines need.  The retained environments print in this document as Corollary 1, Proposition 1 in order of appearance, whereas the article prints the same items with the numbering of its own full document.  The complete native source of the article as distributed in the arXiv e-print for this exact version is bundled unchanged as \texttt{sources/192157/source0.tex}.
\begin{corollary}
    \label{cor:gen_fun}
    Let $\GG_{-n}$ be the set of isomorphism classes of half-edge labelled graphs with Euler characteristic $-n$. Then the following identity holds true.
    \begin{equation}
    \label{eq:gen_fun_euler}
    \sum_{\substack{n\geq 0\\ G \in \GG_{-n}}} \frac{z^{\chi(G)}}{|\Aut(G)|}\prod_{v\in V_G} \Lambda_{\deg(v)} =  \sum_{\bb{s} \geq 0} z^{|\bb{s}|} (2s_1-1)!!\cdots (2s_d-1)!! [\bb{x}^{2\bb{s}}]\exp\left( \sum_{\substack{\bb{w}\in \ZZ^d_{\geq 0}\\ |\bb{w}| \geq 1}} z\Lambda_{\bb{w}} \frac{\bb{x}^{\bb{w}}}{\bb{w}!} \right)
  \end{equation}  
\end{corollary}

\begin{proposition}
  \label{prop:An_exp_integral}
  Let $V\in \QQ[\lambda][\bb{x}]$ be homogeneous of degree $k$ in $\bb{x}$, and let us denote $M = \tfrac{k}{k-2}$ and $K = \tfrac{2}{k-2}$. Then, for any integer $n\geq 0$ such that $nK,nM\in\ZZ$, we have
  \begin{equation}
    \label{eq:An_exp_integral}
    A^V_n (\lambda) = \frac{2^{nM + (d-2)/2} \left(nM +\frac{d-2}{2}\right)!}{(2\pi)^{d/2} (nK)!} \int_{S^{d-1}} V(\bb{x},\lambda)^{nK} \,\omega_{S^{d-1}},
  \end{equation}
  where $\omega_{S^{d-1}}$ denotes the standard volume form of the $(d-1)$-dimensional sphere $S^{d-1}$.
\end{proposition}

\begin{proof}
  Since $|V_G|$ and $|E_G|$ must be integers, $A^V_n = 0$ unless $nK,nM\in\ZZ$. So let us assume the latter. From Corollary \ref{cor:gen_fun} we have 
  \begin{equation}
    \label{eq:An_series}
    A^V_n(\lambda) = \sum_{\substack{\bb{s} \geq 0 \\ |\bb{s}| = nM}} (2s_1-1)!!\cdots (2s_d-1)!!\cdot [\bb{x}^{2\bb{s}}]\exp\left( \sum_{\bb{w}\in \ZZ^d_{\geq 0},\, |\bb{w}| \geq 1} \Lambda_{\bb{w}}(\lambda) \frac{\bb{x}^{\bb{w}}}{\bb{w}!} \right).
  \end{equation}
  The $\bb{x}^{2\bb{s}}$ coefficient of the exponential on the left-hand side has degree $\tfrac{2nM}{k} = nK$ in the coefficients $\lambda_{\bb{w}}$. Using the homogeneity of $V$, this implies
  \[
    [\bb{x}^{2\bb{s}}]\exp\left( \sum_{\bb{w}\in \ZZ^d_{\geq 0},\, |\bb{w}| \geq 1} \Lambda_{\bb{w}}(\lambda) \frac{\bb{x}^{\bb{w}}}{\bb{w}!} \right) = [\bb{x}^{2\bb{s}}] \frac{V(\bb{x},\lambda)^{nK}}{(nK)!}.
  \]
  Recall the Gaussian integral identity
  \[
    \frac{1}{\sqrt{2\pi}} \int_{-\infty}^{\infty} e^{-\frac{x^2}{2}} x^t \d x = \left\{ \begin{array}{ll}
      (t-1)!! & \text{if } t \text{ is even,} \\
      0 & \text{if } t \text{ is odd.}
    \end{array} \right.
  \]
  This enables us to rewrite \eqref{eq:An_series} as 
  \[
    A^V_n(\lambda) = \frac{1}{(2\pi)^{d/2} (nK)!} \int_{\RR^d} \exp\left( -\frac{1}{2} \sum_{i=1}^{d} x_i^2 \right) V(\bb{x},\lambda)^{nK} \d\bb{x}.
  \]
  Finally, the homogeneity of $V$ allows rescaling $\bb{x}$ to have unit norm via $V(r\cdot \bb{x},\lambda) = r^k V(\bb{x},\lambda)$, so we can rewrite the integral above in polar coordinates as
  \begin{align*}
    A^V_n(\lambda) & = \frac{1}{(2\pi)^{d/2} (nK)!} \int_{0}^{\infty} r^{d-1} \left( \int_{\partial B_r(0)} \exp\left( -\frac{1}{2} \sum_{i=1}^{d} x_i^2 \right) V(\bb{x},\lambda)^{nK} \,\omega_{\partial B_r(0)} \right) \d r\\
    & = \frac{1}{(2\pi)^{d/2} (nK)!} \int_{0}^{\infty} e^{-\frac{r^2}{2}} r^{d-1 + knK} \d r \int_{S^{d-1}} V(\bb{x},\lambda)^{nK} \,\omega_{S^{d-1}}.
  \end{align*}
  Together with the identity
  \[
    \int_{0}^{\infty} e^{-\frac{r^2}{2}} r^{nkK+d-1} \d r = \int_{0}^{\infty} e^{-q} (2q)^{(knK + d-2)/2} \d q = 2^{nM + (d-2)/2} \left(nM + \frac{d-2}{2}\right)!\,,
  \]
  this yields the desired expression.
\end{proof}

\end{document}
